\begin{theorem}[sum-moebius-sq-divisors-apply-prime-pow]
\uses{sum_moebius_sq_divisors}
\label{sum_moebius_sq_divisors_apply_prime_pow}
\lean{ArithmeticFunction.sum_moebius_sq_divisors_apply_prime_pow}
% at /home/runner/work/trough/trough/library/PrimeNumberTheoremAnd/IwaniecKowalskiCh1.lean:1683.0-1719.67
\leanok
Applied at prime powers, sum-moebius-sq-divisors coincides with $\mu^2$.
\end{theorem}
\begin{proof}
\uses{sum_moebius_sq_divisors_apply}
\leanok
For a prime power $p^k$, note that if $k\leq 1$ then the only square divisor is $1$, so the sum
    evaluates as $\mu(1)=1$. If $k\geq 2$, then $1$ and $p^2$ are square divisors of $p^k$. Thus,
    the sum evaluates as $\mu(1)+\mu(p)+\ldots$ where the remaining terms are moebius of higher
    powers (if necessary). Since $\mu(p)=-1$ and $\mu$ of higher powers of $p$ is zero, this is
    $0$. This agrees with $\mu(p^k)^2$, which is simply an indicator function for $k\leq 1$ (i.e.
    $p^k$ is squarefree).
\end{proof}
